% Portable theorem guard: define only what the host preamble is missing, so this
% file drops into any of our manuscripts and re-running it is a no-op.
\makeatletter
\@ifundefined{theorem}{\newtheorem{theorem}{Theorem}}{}
\@ifundefined{lemma}{\newtheorem{lemma}[theorem]{Lemma}}{}
\@ifundefined{proposition}{\newtheorem{proposition}[theorem]{Proposition}}{}
\@ifundefined{corollary}{\newtheorem{corollary}[theorem]{Corollary}}{}
\@ifundefined{definition}{\newtheorem{definition}[theorem]{Definition}}{}
\@ifundefined{remark}{\newtheorem{remark}[theorem]{Remark}}{}
\makeatother
% BT1175 triple-45 bridge theorem.

\begin{theorem}[Triple model for the $45$-point incidence object]
\label{thm:triple-fortyfive-incidence-bridge}
Under labeled identifications, the following three $45$-point incidence objects
carry the same point graph parameters:
\[
  (v,k,\lambda,\mu)=(45,12,3,3):
\]
\begin{enumerate}
\item the Boolean/Clifford relation $45$ from the orbit-closed feature module;
\item the cubic-surface tritangent $45$ in the three-layer model;
\item the $W(3,3)$ center-quad quotient point graph.
\end{enumerate}
The transport graph previously constructed from the center-quad quotient is the
complement, with parameters $(45,32,22,24)$.  Thus the Boolean relation sector,
the tritangent layer model, and the W33 center-quad quotient are the same labeled
incidence object at the point-graph level.  The construction remains labeled:
full label-free $S_6\simeq Sp(4,2)$ naturality of the mask dictionary is a
separate refinement.
\end{theorem}
